\section{Related Work}
\label{sec:related}

\para{Host-interconnect contention and congestion control.}
Neugebauer \etal{}~\cite{neugebauer2018pcie} characterized PCIe transaction bottlenecks in end-host networking, while Host Congestion Control~\cite{agarwal2023hostcc}, Hostping~\cite{liu2023hostping}, RHCC~\cite{wan2025rhcc}, and \emph{Understanding the Host Network}~\cite{vuppalapati2024hostnet} demonstrated that high-rate \emph{inbound} NIC-to-host DMA saturates host memory controllers and I/O links, motivating host-aware congestion control. Our work uncovers the symmetric \emph{outbound} pathology on modern x86 servers: unpaced host-to-device posted MMIO writes (\texttt{MOVDIR64B}) bypass the LLC and DDR5 channels yet exhaust the shared 32-slice \texttt{CHA TOR} once \texttt{WCiLF} crosses 330--430 entries, starving co-located compute even when 22--25\,GB/s of DDR5 headroom sits idle and even on an on-die 4\,KiB accelerator portal with zero DMA.

\para{PIO, MMIO direct push, and streaming NIC interfaces.}
High-performance RDMA and RPC stacks (FaSST~\cite{kalia2016fasst}, RDMA design guidelines~\cite{kalia2016rdma}, eRPC~\cite{kalia2019erpc}, CliqueMap~\cite{singhvi2021cliquemap}) long exploited BlueFlame/PIO to cut small-message TX latency. More recently, Liew \etal{}~\cite{liew2025hotos,liew2026asplos} argued that MMIO is the shortest path between CPU and NIC and proposed hardware remote-memory ordering extensions, while Ruzhanskaia \etal{}~\cite{ruzhanskaia2024pio} revisited PIO on coherent interconnects (Enzian~\cite{cock2022enzian}) and Ens\={o}~\cite{sadok2023enso} redesigned the NIC--CPU streaming abstraction to amortize PCIe crossings. Our measurements complement these designs on off-the-shelf PCIe Gen5 and on-die DSA silicon: we quantify the exact pre-\texttt{IIO} 32-\texttt{CHA} reordering rates, show why naive multi-core MMIO push triggers socket-wide \texttt{CHA TOR} collapse, and prove that pacing plus frontier-safe device consumption unlocks direct push safely on today's hardware.

\para{Cache-coherent NICs and DDIO/RX optimization.}
CC-NIC~\cite{schuh2024ccnic}, CXL-NIC~\cite{cxlnic2025}, Dagger~\cite{lazarev2021dagger}, and nanoPU~\cite{ibanez2021nanopu} explore cache-coherent or CPU-integrated NIC attachments. Meanwhile, ResQ~\cite{tootoonchian2018resq}, Farshin \etal{}~\cite{farshin2020ddio}, ShRing~\cite{pismenny2023shring}, rxBisect~\cite{pismenny2025rxbisect}, CEIO~\cite{liu2025ceio}, Sumeru~\cite{wang2026sumeru}, and Sepia~\cite{song2026sepia} optimize \emph{NIC-to-host RX} cache utilization and DDIO~\cite{intelddio} isolation. Our paper focuses squarely on the \emph{host-to-device TX} direct-push path on deployed non-coherent PCIe Gen5 and on-die accelerator silicon, avoiding speculative CXL hardware assumptions and pairing naturally with RX cache-isolation designs.

\para{Lock-free queues and SmartNIC internal contention.}
Classic SPSC ring designs (Lamport~\cite{lamport1977concurrent,lamport1983specifying}, FastForward~\cite{giacomoni2008fastforward}, MCRingBuffer~\cite{lee2010mcringbuffer}, B-Queue~\cite{wang2013bqueue}, Maffione \etal{}~\cite{maffione2019spsc}) and spin-wait algorithms~\cite{anderson1990spinlock,mellorcrummey1991mcs,karlin1991competitive} optimize symmetric CPU-to-CPU cache coherence. We show that asymmetric PCIe-to-SmartNIC rings invert key assumptions---most strikingly, B-Queue's speculative probing plants ``snoop mines'' that trigger a 2.5$\times$ collapse on CN106 \texttt{Normal-WB} pages. Finally, Wei \etal{}~\cite{wei2023smartnic} and Liu \etal{}~\cite{liu2026smartnic} characterized SoC SmartNIC internal bottlenecks; we drill into the exact CHI/TAD coherence boundary (\texttt{Normal-WB} 13.4\,M/s snoop lock vs. zero-\texttt{MAIR} \texttt{MT\_NORMAL\_NC} 174\,Mpps and \texttt{EXLMN} spillover) linking device-internal contention directly back to host \texttt{CHA TOR} starvation.
